\chapter{Conclusion}
\label{ch:conclusion}

\section{What was built}

HALQA is a working full-stack GenAI peer-learning companion, delivered as a NestJS 11 API
over MongoDB with a React 19 progressive web app. Its distinguishing property is not the
model integration but the discipline around it: a documented two-tier routing policy that is
measured rather than asserted, a four-state degradation machine in which every state changes
what the interface promises, a retrieval layer whose citations are mandatory and reusable
for verification, and a privacy guarantee expressed as a build gate.

The product surface is real. Twenty-five backend modules and sixty-eight routes support
timed peer windows, a learning-DNA passport, a leak-gated hint ladder, Auditor verification
with poll-first fallback, consensus polls, revision capsules, quiz battles, the Whisper
chain-repair game, S-B-U circle health, a deterministic matcher, a mentor seat, solo mode,
a facilitator board and an observability surface. The client ships fifteen routes in four
registers with right-to-left support and an offline pack.

\section{What was measured}

Against the real inference platform, a 45-check smoke walks every endpoint in dependency
order and completes with zero failures, reporting a tier split of C = 13, L = 2, that is
\code{lShare = 0.133} against a $<0.25$ ceiling. Two hundred and twelve automated
assertions across three tiers plus five replayable evaluation sets back the guarantees the
report claims. The peer-state poll has a p95 objective treated as a pass/fail gate, not as a
dashboard.

\section{What was learned}

Four lessons generalize beyond this project, and they are the ones worth carrying forward.

\paragraph{State machines beat instructions.} Every behaviour that mattered --- withholding
the AI answer until the window closes, locking the fourth hint level, refusing a quota
reservation --- is enforced on persisted state rather than written into a prompt. A prompt
instruction is a suggestion that a determined input can violate; a condition on state is a
constraint.

\paragraph{Measured policy survives review; asserted policy does not.} The two-tier split
was designed to be falsifiable: the share is computed from the call ledger and published as
a product metric with a target. The same code would look identical if the policy were bad,
which is exactly why the number had to exist.

\paragraph{Failure design is part of the product.} The degradation machine took more design
attention than the happy path, and it shows in the result: the circle keeps working when the
model does not, and the interface says what it is no longer promising. Systems that only
design for availability discover their failure story in front of users.

\paragraph{Privacy is cheap when it is structural.} Removing the columns in which personal
data could live cost nothing at design time and makes the guarantee testable forever. The
equivalent privacy policy, written in prose, would have cost a page and proved nothing.

\section{Perspectives}

\paragraph{Evaluation with real cohorts.} The gap between a passing smoke and a better
learning outcome is the central open question. Real-user testing in two contexts --- a
university tutorial and a self-learning group --- is planned for the post-submission rehearsal,
with the peer-first rate, the S-B-U rings and the confusion signals recorded alongside the
four-language interaction evidence.

\paragraph{Tier quality comparison.} The routing policy now needs the experiment it implies:
the same question set answered at both tiers and scored for pedagogical quality, so the
escalation gate can be tuned against outcome rather than against cost.

\paragraph{Generalization of the pack format.} The knowledge base is deliberately a data
format, not a model, so the interesting test of the design is how cheap it is for an
adopter to add a course. If a new pack takes an afternoon, the approach scales; if it takes
a sprint, it does not.

\begin{keyfinding}[Closing]
The AI does not replace the circle --- it makes the circle work. On the evidence in this
report, the mechanism for that is implemented, measured and defensible; the claim that it
improves learning is the next experiment, and the report is explicit about the difference.
\end{keyfinding}
